%───────────────────────────────────────────────────────────────────────────────
\subsection{JVP Strategies}
\label{sec:jvp_strat}

Two strategies are implemented for computing the directional derivative in
Eq.~\eqref{eq:jvp_def}.

\textbf{Strategy 1: Exact JVP} (\code{torch.func.jvp}).  Computes the exact
directional derivative via forward-mode automatic differentiation:
\begin{equation}
  \bigl(u_\theta,\;\tfrac{du}{dt},\;\hat{v}\bigr)
    = \texttt{torch.func.jvp}\!\bigl(
        \underbrace{f_{\mathrm{dual}}}_{\mathbb{R}^d \times \mathbb{R} \times \mathbb{R} \,\to\, \mathbb{R}^d \times \mathbb{R}^d},\;
        \underbrace{(z_t,\,t,\,r)}_{\text{primals}},\;
        \underbrace{(\tilde{v},\,\mathbf{1},\,\mathbf{0})}_{\text{tangents}},\;
        \texttt{has\_aux}{=}\texttt{True}\bigr)
  \label{eq:exact_jvp}
\end{equation}
\begin{shapes}
Primals: $z_t \!:\! (B, 4, 48, 48, 48)$; $t \!:\! (B,)$; $r \!:\! (B,)$.\\
Tangents: $\tilde{v} \!:\! (B, 4, 48, 48, 48)$; $\mathbf{1}_B \!:\! (B,)$; $\mathbf{0}_B \!:\! (B,)$.\\
Returns: $u_\theta \!:\! (B, 4, 48, 48, 48)$; $du/dt \!:\! (B, 4, 48, 48, 48)$; $\hat{v} \!:\! (B, 4, 48, 48, 48)$.\\
\code{has\_aux=True} captures $\hat{v}$ without computing its JVP (zero overhead).\\
Compound: $V = u + (t-r) \cdot \sg[du/dt]$: $(B, 4, 48, 48, 48)$.\\
\textbf{Code:} \code{jvp\_strategies.py:109--119}.
Memory: ${\sim}20$\,GB activations per sample at $B{=}2$ on A100-40GB.
\end{shapes}

\noindent
Here $f_{\mathrm{dual}}(z,t,r) \coloneqq (u_\theta(z,t,r),\;\hat{v}(z,t,r))$
returns both the average velocity (through which the JVP is computed) and the
v-head output (captured as auxiliary, no JVP).  The function
\code{\_u\_with\_v\_aux} wraps the dual model output so that
\code{has\_aux=True} passes the v-head through unchanged.

\textbf{Strategy 2: Finite Difference JVP} (FD-JVP).  Approximates the
directional derivative using a perturbation of step size
$\delta = 10^{-3}$:
\begin{equation}
  \frac{du_\theta}{dt}\bigg|_{\tilde{v}}
    \approx \frac{
      u_\theta^{\mathrm{(fp32)}}(z_t + \delta\,\tilde{v},\;t + \delta,\;r)
      - u_\theta^{\mathrm{(fp32)}}(z_t,\;t,\;r)
    }{\delta}
    \;\in \mathbb{R}^{C \times D \times H \times W}
  \label{eq:fd_jvp}
\end{equation}
\begin{shapes}
Perturbed inputs: $z_t + \delta\tilde{v}$: $(B, 4, 48, 48, 48)$; $t + \delta$: $(B,)$.\\
Both $u$ terms cast to fp32 before subtraction (avoids bf16 catastrophic cancellation).\\
Perturbed pass runs under \code{torch.no\_grad()} (since $du/dt$ is always detached).\\
\textbf{Code:} \code{jvp\_strategies.py:156--168}.
\end{shapes}

\textbf{Critical incompatibility (Phase~4f).}  x-prediction + FD-JVP is
\textbf{numerically unstable}.  The x-to-u conversion
$u = (z_t - \hat{x})/t$ has a $1/t$ singularity.  When FD-JVP computes
the difference quotient, the result involves:
\begin{equation}
  \frac{du}{dt}\bigg|_\mathrm{FD}
    = \frac{u(t{+}\delta)-u(t)}{\delta}
    \sim \frac{O(1/t)}{\delta}
    = O\!\left(\frac{1}{t^2\,\delta}\right) \to \infty
  \quad\text{as }t\to t_{\min}
  \label{eq:fd_instability}
\end{equation}
With exact JVP, the $1/t$ factor is analytically differentiated via the
chain rule, yielding stable gradients.

\textbf{Rule:} \fbox{x-pred + exact JVP = stable}\ \ %
\fbox{u-pred + FD-JVP = stable}\ \ %
\fbox{x-pred + FD-JVP = explosion}
